\par\Needspace{10\baselineskip}\pdfbookmark[1]{Matthiessen. 1819 January 30.}{matthiessen1819}\medskip\hrule\smallskip
\begin{center}
Göttingische gelehrte Anzeigen. 1819 January 30.
\end{center}
\smallskip\hrule\medskip

[E. A. \textsc{Matthiessen}.] \textit{Tafel zur bequemern Berechnung des Logarithmen der
Summe oder Differenz zweier Grössen, welche selbst nur durch ihre Logarithmen gegeben sind.} Altona 1818. Published by J. F. Hammerich. Introduction, 33 pp.; tables,
212 pp., quarto. (Title and introduction also in Latin.)

In somewhat lengthy calculations, every shortcut is valuable; even aids that offer only a small advantage when used once become important through repeated use. One such aid is an auxiliary table first given in the Monatliche Correspondenz in 1812, for obtaining the logarithms of the sum or difference of two quantities directly from their logarithms: a single lookup achieves what would otherwise require three, or at least two. With such a table, one can usefully retain many \SourceLeaf{260} formulas in their original form, instead of seeking to make them more convenient for calculation by introducing auxiliary angles. The table mentioned was intended only for calculations with five decimal places: precisely because such calculations occur most frequently, this was the most immediate need. The wish expressed when it was published, that someone might take the trouble to calculate a similar table on a larger scale and with seven decimal places, led to the present tables. Through this laborious work their author, Mr \textsc{Matthiessen}, has earned the gratitude of everyone who does much logarithmic calculation. The preface describes the method used by the author in calculating the table, and his commendable care to make it reliable throughout, even in the last digit. In tables of this kind, the outward arrangement also matters. The author initially had an arrangement in mind that would compress the table into the smallest possible space. But since the table's entire purpose is to \textit{make calculations easier}, that purpose would be completely defeated if its arrangement were too intricate and required burdensome attention. On further consideration, the author therefore rightly decided to set aside that ingenious arrangement and choose a simpler one instead, although this considerably increased the size of the volume. Perhaps some users of the tables will join us in wishing that the author had gone a little further in this respect and printed each of the columns \(B\) and \(C\) in full, rather than printing only once their last four digits, which are common to both. This would have made the table still easier to take in, without increasing its format if slightly smaller type had been chosen. That could have been done without disadvantage, and perhaps even to the benefit of its appearance. Convenience also suffers somewhat from the almost excessive economy with digits, where as a rule only the fourth is printed while the preceding ones remain unchanged, and with proportional parts, which are always given only once and thus sometimes require the user to turn backwards or forwards. Finally, we would have wished the last part of the table, p. 212, to have been given over a range ten times as large, to make the second differences imperceptible; this would have required \SourceLeaf{261} only one additional page. All these remarks, which may in part depend on individual habits, are in no way intended to diminish the merit of this work, which will certainly be acknowledged by everyone who has occasion to use it.

